\documentclass[11pt,a4paper,twoside]{article}
\usepackage[T1]{fontenc}
\usepackage[inner=22mm,outer=68mm,bindingoffset=3mm,top=30mm,bottom=24mm]{geometry}
\usepackage[grown,side=outer,style=quiet,pagenumber,number-scale=0.85,
  variation=spread,seed=7,color-mode=mono,sections]{florilegium}
\usetikzlibrary{calc}
\pagestyle{empty}
\setlength\parskip{0.4em}
\pagenumbering{roman}\setcounter{page}{14}
\begin{document}
\section{A quieter page}
A book needs rhythm as well as ornament. Here the garden occupies the outer
margin, leaving the binding and the text clear. The next page uses the same
spread seed, so the two gardens belong together without repeating their page labels.

The quiet treatment keeps one principal flower and leaves more paper between
the smaller motifs. Monochrome converts the illustration's blue washes and
copper accents to gray while retaining their tonal differences. It is useful
for a black-and-white edition or a print proof.

The Roman page label is drawn as planted strokes, with the same curls and small
leaves as the garden. It remains real, searchable text in the PDF. Its lettering
will stay upright when the drawing moves to the other side of the book.

\florilegiumdivider

The note in the margin has a physical reservation around it. Stems take a detour
where space permits; motifs that would overlap it are omitted. A final clip
keeps its paper clear. This is a selected reservation, not an automatic rule
for every figure in the document.

\florilegiumreserve[3mm]{10mm}{108mm}{31mm}{36mm}
\AddToShipoutPictureFG*{\begin{tikzpicture}[remember picture,overlay]
  \node[anchor=south west,draw=black!30,fill=black!3,inner sep=2mm,
    text width=27mm,font=\footnotesize] at ($(current page.south west)+(10mm,108mm)$)
    {An intentional note. The plant leaves this part of the margin clear.};
\end{tikzpicture}}
\clearpage

\section{The facing page}
The seed is shared by pages fourteen and fifteen. The available space and the
text marks still matter: a reservation or a heading can change what grows
nearby. Spread variation provides a common composition, rather than freezing
every leaf regardless of the page's contents.

The ornament is reflected onto the right, but the page label is constructed in
its normal reading direction. Search, selection and copy use the original
lowercase Roman label. The ordinary footer has been turned off explicitly in
this example to leave one visible page number.

\begin{florilegiumprotect}[3mm]
\hrule\smallskip
\small A protected box travels with its contents, including when it is placed
inside a figure or table. It is unbreakable; use separate boxes for material
that spans several pages.
\smallskip\hrule
\end{florilegiumprotect}

\florilegiumdivider[28mm]

The divider is made from the same leaves and rosette as the margin. Its two
halves replay one seed, giving a small, balanced mark between sections.
It is vector artwork and can also be exported as a separate PDF or SVG.

These pages can be built with either LuaLaTeX or XeLaTeX. The drawing and layout
code is shared. The build command manages the extra vector-rendering pass
needed by XeLaTeX and repeats the document until its records and references settle.
\clearpage

\pagenumbering{arabic}\setcounter{page}{38}
\florilegiumsetup{style=chapter,color-mode=color,number-scale=1}
\section{A chapter in blue}
The chapter treatment returns to color, with two principal flowers and a larger
focal scale. This is the same drawing system as the quieter opening pages:
composition, density and palette are separate choices.

A warm accent distinguishes the heading's flower and the fine header line.
The blue ramp keeps the pale washes related to the dark stems. The page label
grows from the baseline at the foot of the margin rather than sitting beside
the plant as a separate badge.

\florilegiumdivider

For a facing pair, keep the settings unchanged between the two pages and use
\texttt{variation=spread}. For a repeated design, choose \texttt{variation=fixed}.
A numbered seed makes either choice reproducible. Text marks, seasons and
reservations can still give each page a different job to do.

\begin{verbatim}
\florilegiumsetup{
  style=chapter,
  variation=spread,
  seed=7,
  pagenumber=true
}
\end{verbatim}
Apply a style change after \verb|\clearpage|: settings are read when the page
is shipped out. This keeps the preceding page's treatment intact.
\clearpage

\section{A coordinated spread}
The page label changes to thirty-nine, while the spread keeps its principal
composition. The lettering reads from left to right on both sides. Its stroke
width follows a broad nib, and the small leaves connect it to the rest of the garden.

There is room for more than one kind of book here. The quiet monochrome pages
would suit a short essay; the richer pair could introduce a chapter, a letter
or a collection of field notes. The controls let the author choose that degree
of ornament without changing the prose.

\florilegiumdivider[32mm]

Reusable motifs can be exported independently through the native PDF renderer.
The same PDF included several times is embedded once and reused. SVG export
preserves vector paths for applications outside TeX.

Every addition in this guide is optional. The original color composition,
page variation and ordinary footer remain the defaults. Select the page-number
option when there is enough margin space, and retain a normal footer for labels
or layouts that the botanical lettering cannot safely accommodate.

\smallskip\noindent\textit{Built from curves, placed around the text.}
\end{document}
